\section{规范中的五处张力}
\label{sec:synthesis}

对照证据阅读，该规范在五处向两个方向拉扯，而低成本的类型化判断使其中每一处都从一个哲学问题变成一项工程默认设置。
对每一处张力，我们给出对文本最善意的解读、与之相关的证据以及若干选项，其中至少有一个选项保留原文不变。
在各选项之间作出选择，属于该框架的作者以及将来实施它的人。

\paragraph{T1. 理解语言中的情感与检测个人的情感状态。}
\Clause{4.3} 指出，“毫无必要让 PAI 检测一个人的情感状态”，因为情感状态是转瞬即逝的、个性化的、情境敏感的，并涉及人的尊严；\clause{4.3.3.1} 要求 PAI “理解人类情绪状态和过程”；\clause{4.3.3.3} 则以两个关于小狗狗的问句说明语义识别，二者的区别在于话语的意图。
将这些条款合读，可以得到一种一致的解释：识别话语的情感意义与意图，但绝不推断个人的内心状态。
然而，文本并未明示这一边界，而类型化模型使得为每条消息附加情绪分数几乎毫无成本，这正是该边界在实践中被越过的方式。
关于情感相关判断的证据薄弱且结果不一：在一项标注研究中，\jev{} 对大多数共情条目给出了高置信度，但在其置信度最高的条目上，准确率几乎不高于全部条目（0.383 对 0.371，三类均衡），作者在一个低成本生成式模型中也发现了同样的失效~\cite{arx24574}；而在多语言效价--唤醒度回归任务上，加装了一个针对标签拟合的监督层的 \jev{} 与最佳系统持平~\cite{arx35293}。
更广泛的文献则更为明确：从可观察信号推断情绪在科学上依据薄弱~\cite{barrett2019emotional}，在伦理上问题重重~\cite{stark2021ethics}。
\emph{选项：}明确写出这一边界，使保障机制判断话语的语义与意图，但绝不为个人赋予情感状态；并将情感相关判断保留给 PAI 审计其\emph{自身}输出之用，例如审计 \clause{4.3.2} 所称构成情感操控的虚假亲密关系。

\paragraph{T2. 二值判定与三值伦理。}
\pol{} 建立在爱/恨两极之上，但增设了不在场状态（\clause{4.3.2}），并告诫其概念服务于结构性分析而非道德标签（\clause{1.5}）。
“这是恨语吗？”这样的是/否问题把这两项限定都舍弃了。
证据从多个方向指向同一结论：选项名称可以左右判定结果~\cite{arx26758,arx00346}；在模型不确定之处，互补概率之和不等于一~\cite{arx33209}；显式的“不知道”选项在其为正确答案时很少被选中，“以上皆非”也只是不稳定地被选中（在一个医学基准上为 54\%，在识别它需要算术推理时为 7\%）~\cite{arx34024,arx39496}；是/否概率把冲突与无知置于同一区间，而一个分别列出这两种状态的 Choice 问题在四个案例中把二者区分开来~\cite{zen22848952}\zmark{}。
证据也表明了什么有帮助：单独询问证据是否充分的是/否问题，在发现信息缺失方面优于置信度，尽管它作为接受答案的筛选器较弱~\cite{arx01006}。
\emph{选项：}将每个爱/恨问题都设为三值，把不在场或信息不足列为显式选项，并将充分性作为单独的问题提出、监测其使用率；将“有争议”（人们意见分歧）作为一个单独的标记交由人处理，与“不在场”相区分；并通过把中性标识符绑定到定义、再以重命名加以检验的方式，使道德词汇不出现在选项名称中。
这些选项是在实施 \clause{1.5} 与 \clause{4.3.2}，而非改变它们。

\paragraph{T3. 完全透明与封闭模型。}
\Art{4} 与 \art{5}(4) 要求运行逻辑和决策过程“完全公开透明”，并要求技术“不应是黑箱”。
TypeSafe 既未公开托管版 \jev{} 的权重、架构、奖励，也未公开其训练数据，且不提供自托管~\cite{arx24574,arx28919,arx00346}。
一种善意的解读是，\art{4} 关注的是 NaturalDAO 的流程与资源流向的透明，而记录一个封闭组件的输入、输出、版本与校准数据可以部分满足这一要求；但此类日志无法显示该组件如何得出其输出，而仅凭开放权重也无法提供 \art{5}(5) 所要求的“推理步骤”。
开源模型使读出过程可被检查，但在若干比较中其准确性或鲁棒性较差：Laya 表现出最严重的名称反转，在交互风险筛查上接近随机水平~\cite{arx26758,arx33401}。
这一差距并非固定不变：更大的开源解码器在选项名称互换下比 \jev{} 更鲁棒~\cite{arx00346}；由 Laya 初始化的中文模型在其作者的通用基准上、以及经领域微调后在医学领域略微超过 \jev{}，尽管在各专业领域上的平均准确率为 \jev{} 的 92\%~\cite{arx36965}；而且已有一套开放的 RLCD 方案~\cite{arx38850}。
\emph{选项：}（a）仅将封闭模型用作外部基线，而在可自托管的模型上运行具有实质后果的判断，并在每条决策记录中固定其权重、版本与校准数据；或（b）允许封闭模型作为若干独立检测器之一，同时记录其输入、输出与版本。
选项（a）以已测得的准确性与鲁棒性换取可检查性；这一权衡应明确作出，并随开源模型的进步而重新审视。
公开决策记录还与不向攻击者泄露概率相冲突（见\cref{sec:g2}），\cref{sec:design} 通过延迟并粗化公开内容来化解这一冲突。

\paragraph{T4. 纯 AI 裁决与相关错误。}
\Art{2}(1) 规定 NaturalDAO “在与全人类密切协作的前提下”自主完成决策，\art{10} 将争议交由 AI 进行事实分析与伦理评估，并接受人类质询，\art{9}(7) 则规定人类监督“不直接干预”决策，质询中提出的问题由 NaturalDAO 自主处理。
文本并未说明协作这一前提如何约束具体决策。
证据从两方面涉及非干预原则。
升级可能保留错误：在基于评分细则的评分任务上，低成本生成式评估器在 96\% 的情形中重复了 \jev{} 置信度最高的错误（这是一个上界，因为条目是作为高置信错误被挑选出来的），而在智能体安全筛查中，评估器至多发现了 130 个高置信漏检中的 5 个~\cite{arx29769,arx33401}。
但当下一阶段是更强的推理模型或独立构建的规则时，升级也可以奏效~\cite{arx26550,arx02048}，尽管后来的一项研究使前一种情形有争议（\cref{sec:update}）。
人类作为终点在治理条目上能否做得更好尚未经检验：没有研究在此类条目上比较人类审核者与类型化模型，而监督相关文献表明，人们同样会过度依赖高置信系统，且彼此之间意见不一~\cite{parasuraman2010complacency,green2022flaws,bansal2021whole}。
\emph{选项：}（a）保留非干预原则，但使质询切实有效：公开发布对高置信决策的常设随机审计，规定 NaturalDAO 有义务对每一处差异公开回应，并升级至独立构建的核查器而非同类模型；或（b）通过 \art{11} 的修订程序界定一类福利决策，例如扣发福利（并通过对 \clause{5.4.1} 的相应修改，涵盖阻断某人的言论），人类审核者可在审查期间暂停其自动执行。
选项（a）保留了原文，但将具有实质后果的决策留给自动化组件，因而超出了\cref{sec:intro} 所界定的检测器角色；我们列出它是因为它保留了原文，而非因为证据支持它。选项（b）增加了一道核查，其价值从正当程序文献来看是合理的~\cite{citron2008technological}，且与现行法律相近（见下文），但尚未在该任务上检验；\cref{sec:design} 的参考设计实现了这一选项。

\paragraph{T5. 评估每个人与平等联结。}
在经过治理的 AI “可以”做的事项中，\clause{5.6} 列出了量化“每个人的爱语与恨语特征”，并配合激励措施引导行为。
\Clause{4.3.1} 认为，平等对待每一个人与平等对待每一个人的行为“不能划等号”，其上下文也拒绝依据人的所作所为给人贴标签。
类型化模型能以不到一美分的成本为每个人的每句话评分，因此逐人汇总在技术上轻而易举。
这些模型会从默认的文化立场回答价值问题~\cite{arx36399}，在一项非正式的强制选择测试中表现出群体偏见~\cite{simmons2026saidno}\gmark{}，在低资源语言上表现较差（在一个基准中，类型化模型与生成式模型均是如此~\cite{arx37647}），并且可被插入内容中的第三方观点所操纵~\cite{arx31142}。
若对一个人的历史进行汇总，此类错误可能不均衡地落在不同群体身上，并成为一个可被他人操纵的持久标签；目前尚无研究测量行为判断中按说话者群体区分的差异化错误。
社会评分系统的经验表明，此类评分如何成为控制的基础设施~\cite{creemers2018social,liang2018constructing}。
\emph{选项：}对事件与行为评分，绝不对人评分，且不保留任何持久的逐人汇总；或者，若实施 \clause{5.6}，则对任何逐人度量要求取得同意、允许本人查阅自己的记录、提供申诉途径，并按群体与语言公布错误率。

\paragraph{监管上的对应。}
这些张力没有一处是 \pol{} 所独有的，而欧盟法律已直接或以类比方式对其中若干处表明立场。
欧盟《人工智能法》~\cite{euaiact2024} 禁止在工作场所和教育机构中进行情绪识别，出于医疗或安全原因者除外（第 5 条第 1 款（f）项）；由于该法将情绪识别系统定义为依据生物识别数据推断情绪的系统（第 3 条第（39）项；序言第 44 段），仅依据文本进行的推断（即类型化模型所做的推断）不在该禁止范围之内，因此该条款是以类比方式而非直接适用的方式与 T1 相关；它还禁止社会评分，即依据人的社会行为或个人特征对其进行长期评估、且该评分导致其在无关情境中受到不利对待或受到不合理或不相称对待的做法（第 5 条第 1 款（c）项），这与 T5 相关。
对于高风险系统，该法要求有效的人类监督（第 14 条），并赋予受到基于其附件三所列特定高风险系统、具有法律效果或类似重大影响之决策的人，获得关于该系统所起作用及决策主要要素之解释的权利（第 86 条）。这两项规定仅适用于高风险系统，而言论过滤器和私人组织的福利决策均未列入附件三（其第 5（a）点涵盖由公共机构或代表公共机构作出的公共援助决策）；不过，第 8（a）点确实列入了用于替代性争议解决的 AI，当 \art{10} 下的争议裁决结果对当事方产生法律效果时，该点可能适用于此类裁决（序言第 61 段）。因此，就安全阀与福利决策而言，这两项规定是以类比方式与 T4 以及 \art{9}(2) 的解释义务相关，而 GDPR 则在处理个人数据的任何场合均适用。
《通用数据保护条例》（GDPR）限制仅基于自动化处理、产生法律效果或类似重大影响的决策，并在此类决策以合同或明确同意为依据时，至少要求保障获得人工干预、表达本人观点以及对决策提出异议的权利（第 22 条第 1、3 款）~\cite{gdpr2016}，这与 T4 的选项（b）相近。
在内容审核方面，《数字服务法》（DSA）要求托管服务提供者对每一项限制提供理由说明（第 17 条），要求在线平台设立内部投诉处理系统（第 20 条）~\cite{dsa2022}；《圣克拉拉原则》则要求通知、说明理由与申诉~\cite{santaclara2021}；这些都与阻断恨语的安全阀十分相近。
上述选项，即不对个人进行情感状态推断、不保留持久的逐人评分、阻断可由人暂停并附理由与申诉途径、公开决策记录，将使实施更接近这些规范。

\Cref{fig:framework} 显示了每处张力所涉及的要求。
